% Part D. Independent audit and controlled-topology results.
\section{The actual topology of discrete-operator centralizers}
\label{sec:D}

We retain the hypotheses of LM4 C.81, not merely its notation.  Let
$\overline D$ be a closed topological disk, $O=\operatorname{int}\overline D$,
and let $d$ be a compatible metric on $\overline D$.  Let $Q$ be a closed
discrete subgroup of $\operatorname{Homeo}(O)$ acting freely and properly.
The group
\[
 C=\{a\in\operatorname{Homeo}(\overline D):
 a|_{\partial\overline D}=1,\quad aq=qa\text{ on }O\ (q\in Q)\}
\]
has the compact-open group topology, equivalently the uniform topology
on maps and inverses.  We assume that $C$ acts transitively on $O$.
Write $\pi:O\to N=O/Q$, fix $x\in O$, and put $p=\pi x$ and
$F=C_x$.  The quotient is a connected, second countable surface without
boundary.  Its universal cover is the specified $O$; thus it is
aspherical.  None of the arguments below replaces the given
compactification by a chosen hyperbolic compactification.

\statement{D.1}{Proved}{The invariant metric and compact-core collapse}
The function
\[
 d_Q(u,v)=\sup_{q\in Q}d(qu,qv)
\]
is a bounded $Q$-invariant metric inducing the original topology on $O$.
If $A,B\subset O$ are compact, $a_j\in A$, and
$q_ja_j\to z\in\partial\overline D$, then $q_jB\to\{z\}$ uniformly
in the original metric.
\begin{proof}
The group $C$ is a closed subgroup of the Polish group
$\operatorname{Homeo}(\overline D)$: the boundary conditions and each
commutator equation are closed under convergence on compact subsets of
$O$.  Evaluation is a continuous transitive action on the Polish space
$O$.  The open-orbit-map form of Effros' theorem therefore applies;
a proof with precisely these hypotheses is available in
\cite[Theorem 4.1]{DE_Wengenroth2025}.  In particular, if $u_j\to u$ in
$O$, there are $c_j\to1$ in $C$ such that $c_ju=u_j$.  To obtain this
last assertion, choose a nested countable neighborhood basis of $1$,
use openness of its evaluation images, and choose $c_j$ in successively
smaller neighborhoods.  Commutation gives
\[
 d_Q(u_j,u)\leq\sup_{w\in\overline D}d(c_jw,w)\longrightarrow0.
\]
The converse follows from $d\leq d_Q$.  Boundedness, the metric axioms,
and $Q$-invariance follow directly from the definition.

For the collapse assertion, a failure supplies $b_j\in B$ and a
subsequence with $a_j\to a\in A$, $b_j\to b\in B$ and
$d(q_jb_j,z)$ bounded away from zero.  Choose $u_j,v_j\to1$ in $C$
with $u_ja=a_j$, $v_jb=b_j$, and choose $g\in C$ with $ga=b$.
Then
\[
 q_jb_j=v_jg u_j^{-1}(q_ja_j).
\]
The maps $v_jg u_j^{-1}$ converge uniformly on the closed disk to $g$,
and $g(z)=z$.  Hence $q_jb_j\to z$, a contradiction.  This proof
avoids the compact-selection lemma invoked in LM4 C.88.
\end{proof}

\statement{D.2}{Proved}{The exact $C_0$ image and compact-parameter criterion}
For $h\in C_{\operatorname{Homeo}(O)}(Q)$ define
\[
 \delta_h(\pi u)=d_Q(hu,u),\qquad
 V=\{h\in C_{\operatorname{Homeo}(O)}(Q):\delta_h\in C_0(N)\}.
\]
Here $C_0(N)$ consists of continuous functions vanishing uniformly
outside compact subsets.  Restriction $C\to V$ and extension by the
identity on $\partial\overline D$ are inverse isometric topological
group isomorphisms, for the sums of the uniform distances of maps and
inverses.  For a compact metrizable parameter space $P$, an equivariant
family $h_s$ with continuous evaluation $P\times O\to O$ extends
jointly to the closed disk if and only if
\begin{equation}\label{eq:DE-tail}
 \forall\epsilon>0\ \exists K\Subset N:\quad
 \sup_{s\in P,\,\pi u\notin K}d_Q(h_su,u)<\epsilon.
\end{equation}
Joint continuity of the extended inverse family follows as well.
This proves LM4 C.88 with its stated hypotheses.
\begin{proof}
The function $\delta_h$ is continuous by D.1 and constant on each
$Q$-orbit.  If $\bar h$ denotes descent, then
\[
 \delta_{h^{-1}}=\delta_h\circ\bar h^{-1},\qquad
 \delta_{hk}\leq\delta_h\circ\bar k+\delta_k.
\]
Pullback by a homeomorphism of the locally compact space $N$ preserves
$C_0(N)$, so $V$ is a group.

A continuous compact family fixing the original boundary moves every
point outside one compact $A\Subset O$ by less than $\epsilon$,
uniformly in the parameter.  If $\pi u\notin\pi A$, every $qu$ is
outside $A$.  Commutation and the supremum defining $d_Q$ prove
\eqref{eq:DE-tail}.

Conversely suppose \eqref{eq:DE-tail}.  To check continuity at a
boundary point, let $s_j\to s$ and $u_j\to z\in\partial\overline D$.
A subsequence witnessing failure has either $\pi u_j$ escaping every
compact subset of $N$, or all $\pi u_j$ in one compact set $K$.
In the first case \eqref{eq:DE-tail} gives
$d(h_{s_j}u_j,u_j)\to0$.  In the second case finitely many local
covering sections give a compact $A\Subset O$ with $K\subset\pi A$.
Write $u_j=q_ja_j$, $a_j\in A$.  The set
$B=\{h_s(a):s\in P,a\in A\}$ is compact.  D.1 now gives
$h_{s_j}u_j=q_jh_{s_j}a_j\to z$.  These two cases prove continuity.
The resulting map
\[
 P\times\overline D\longrightarrow P\times\overline D,
 \qquad(s,u)\longmapsto(s,h_su)
\]
is a continuous bijection of compact Hausdorff spaces.  Its inverse
is continuous, so no separate unproved inverse-continuity assumption
is needed.  Necessity applied to this inverse family proves its
uniform tail.

Finally commutation gives the exact identity
\[
 \sup_{u\in O}d_Q(hu,ku)
 =\sup_{u\in O,\,q\in Q}d(hqu,kqu)
 =\sup_{w\in\overline D}d(E(h)w,E(k)w).
\]
The same equality holds for the inverses.  Uniform continuity for
$d_Q$ also follows by applying the modulus of continuity of $E(h)$
to every pair $(qu,qv)$, and similarly for $h^{-1}$.
\end{proof}

\statement{D.3}{Proved}{Pointed descent, its image, and its stronger topology}
Let
\[
 I_p=\operatorname{Homeo}(N\mathbin{\mathrm{rel}}p)_0
\]
denote the identity path component in the compact-open topology on
maps and inverses.  Every $a\in F$ descends to $I_p$.  Every $b\in I_p$
has a unique lift $\lambda(b):O\to O$ fixing $x$, and that lift commutes
with every element of $Q$.  Define
\[
 B_p=\{b\in I_p:\delta_{\lambda(b)}\in C_0(N)\}.
\]
Descent gives a group isomorphism $F\to B_p$ and an isometry if $B_p$
is given the metric
\begin{align*}
 d^{\uparrow}(b,c)={}&\sup_{u\in O}d_Q(\lambda(b)u,\lambda(c)u)\\
 &+\sup_{u\in O}d_Q(\lambda(b)^{-1}u,\lambda(c)^{-1}u).
\end{align*}
For compact parameter spaces, continuity in this metric is exactly
compact-open continuity downstairs together with
\eqref{eq:DE-tail} for the actual pointed lifts.  This proves LM4
C.92, including its distinction between the two topologies.
\begin{proof}
A lift $a$ fixing $x$ and commuting with each deck transformation
sends a path from $x$ to $qx$ to another such path.  Therefore its
descent fixes $p$ and induces the identity on $\pi_1(N,p)$, with no
inner-automorphism ambiguity.  Since $N$ is an aspherical surface,
it is based-homotopic to the identity.  One direct construction
chooses a locally finite triangulation with vertex $p$, constructs
the homotopy on its one-skeleton from the equality on $\pi_1$, and
extends across the cell prisms using $\pi_2(N)=0$.  Local finiteness
makes the result continuous; properness of this intermediate
homotopy is not required.

A disk homeomorphism fixing the boundary preserves orientation,
as follows from the boundary map in relative homology
$H_2(\overline D,\partial\overline D)\to H_1(\partial\overline D)$.
Consequently its descent preserves local orientation at $p$, also
when $N$ is nonorientable.  Epstein's based isotopy theorem applies:
$N$ has no boundary; the plane and open-cylinder orientation
exceptions are satisfied, and the open-M\"obius-band exception is
satisfied at $p$.  We use Theorem 6.4(b), its proof on pp.~101--103,
and Appendix A4 for topological homeomorphisms
\cite{DE_Epstein1966}.  The proof supplies a level-preserving
homeomorphism of $N\times[0,1]$, hence continuity of maps and
inverses in the compact-open group topology.  Thus the descent
belongs to $I_p$.

Conversely a pointed isotopy of $b$ to the identity implies
$b_*=1$.  The covering lifting criterion gives a lift fixing $x$;
uniqueness makes it a homeomorphism and makes $b\mapsto\lambda(b)$
a homomorphism.  The equality on $\pi_1$ gives
$\lambda(b)q=q\lambda(b)$ for each actual deck label $q$.
The lifting operation is continuous in the interior compact-open
topology.  To see this without selecting isotopies continuously,
choose a path from $x$ to a given $u$ and lift its $b$-image from $x$.
Path lifting depends continuously on $b$; near $u$, append short
paths in one covering chart.  Finitely many such charts prove
continuity on any compact set, and the same argument applies to
$b^{-1}$.  D.2 now proves the image description and the metric
identity.  Its compact-family criterion proves the final assertion.
\end{proof}

\statement{D.4}{Proved}{Evaluation has a global section in the original topology}
For every quotient $N$ under consideration, compact or noncompact,
there is a continuous map $s:O\to C$ with $s(y)x=y$ and $s(x)=1$.
Consequently
\[
 C\xrightarrow{\ \cong\ }O\times F,\qquad
 a\longmapsto\bigl(ax,s(ax)^{-1}a\bigr)
\]
is a homeomorphism, with inverse $(y,f)\mapsto s(y)f$.
In particular $\pi_j(F,1)\to\pi_j(C,1)$ is an isomorphism for
all $j\geq1$.
\begin{proof}
Choose a smooth structure and a complete Riemannian metric on $N$,
and lift them to $O$.  A smooth plane coordinate on $O$ gives a
smooth contraction $\gamma_y(t)$ from $x$ to $y$, constant when
$y=x$, and depending smoothly on $(t,y)$.  At time $t$ choose on
$N$ a bump vector field centered at $\pi\gamma_y(t)$ whose value
there is $D\pi\,\partial_t\gamma_y(t)$.  This can be done
continuously in $(y,t)$: use the Riemannian exponential map, radial
parallel transport, and a fixed bump function with radius a positive
continuous function smaller than a quarter of the injectivity
radius and smaller than $1$.  The field is smooth in the spatial
variable.  For each fixed $y$ its time-dependent support is compact,
so its flow exists for $0\leq t\leq1$.  On a compact set of values
of $y$, all these supports lie in one compact subset of $N$ by
completeness.  The lifted field is $Q$-invariant, and its flow takes
$x$ along the specified path $\gamma_y$.  Outside the inverse image
of that common compact set, the lifted flow is the identity.

Let $s(y)$ be the time-one lifted flow, extended by the identity on
the original boundary.  Continuous dependence for flows gives
joint interior continuity.  The compact-support instance of D.2
proves continuity for every compact set of parameters $y$ in the
original closed-disk topology; since $O$ is locally compact this
proves continuity of $s$.  When $y=x$ the chosen field is zero, so
$s(x)=1$.  The displayed maps are inverse by direct multiplication;
their continuity follows from continuity of $s$, evaluation, and
the group operations.  The last assertion follows by contracting
the factor $O$.
\end{proof}

\statement{D.5}{Proved}{The compact quotient case is contractible}
If $N$ is compact, $C$ admits a continuous contraction to $1$ fixing
$1$ throughout.  This proves both LM4 C.90 and LM4 C.91, including
the torus and Klein-bottle quotients.
\begin{proof}
For compact $N$, $C_0(N)=C(N)$ and the uniform-tail condition is
vacuous.  D.3 identifies $F$ with $I_p$ also topologically.  For
example, continuity of the inverse follows by applying D.2 to a
convergent sequence and its limit, viewed as a compact parameter
space.  A compact quotient with universal cover $O$ is neither the
sphere nor the projective plane.  The pointed compact-surface
homeomorphism theorem gives contractibility of $I_p$:
\cite[Lemma 2.1(i), pp.~2--3]{DE_Yagasaki2000}.  Its torus and
Klein-bottle exceptions have empty fixed set, so do not apply to
$(N,\{p\})$.

Normalize any contraction $A_t$ by
$J_t(b)=A_t(b)A_t(1)^{-1}$.  After reversing the parameter if
necessary it satisfies $J_0(b)=b$, $J_1(b)=1$, and $J_t(1)=1$.
Transport it to $F$.  If $R_t$ contracts $O$ to $x$ while fixing
$x$, D.4 supplies the explicit contraction
\[
 a\longmapsto s\bigl(R_t(ax)\bigr)
 J_t\bigl(s(ax)^{-1}a\bigr).
\]
All factors lie in the original group $C$, and the formula fixes
$1$.  Hence every positive based homotopy group vanishes as asserted
in LM4 C.90, and the stronger global assertion of LM4 C.91 follows.
\end{proof}

\statement{D.6}{Proved}{Compact-open convergence does not control the lifted tail}
There are operators satisfying all the hypotheses above, with
$N=S^1\times\mathbb R$, and $b_R\in B_p$ converging to $1$ in
$I_p$ while $d^{\uparrow}(b_R,1)>2/5$ for every $R\geq2$.
\begin{proof}
Identify $O$ with $\mathbb R^2$ by
\[
 j(s,y)=\frac{(s,y)}{1+\sqrt{s^2+y^2}},\qquad
 Q=\langle jTj^{-1}\rangle,\quad T(s,y)=(s+1,y).
\]
Use the Euclidean metric on the closed unit disk and
$p=([0],0)$.  This is a free proper closed discrete action.  Its
boundary-fixed centralizer is transitive: a finite sequence of
point-moving isotopies in small disks downstairs lifts to
$T$-equivariant maps of uniformly bounded Euclidean displacement,
which extend by the identity under $j$; choose the projected path
so that its lift joins any prescribed pair of points upstairs.

Put $\beta(v)=\max(0,1-4|v|)$ and
\[
 A_{R,t}(s,y)=(s+tR\beta(y-R),y),\qquad0\leq t\leq1.
\]
These are equivariant homeomorphisms, have the displayed inverse
with $t$ replaced by $-t$, and fix $(0,0)$.  Their displacements
are bounded in the plane, so their conjugates by $j$ extend by
the identity to the closed disk.  Their descents are supported
in $S^1\times[R-1/4,R+1/4]$.  Thus $b_R=\overline A_{R,1}$ and
$b_R^{-1}$ converge to $1$ uniformly on every compact subset of $N$.
Their unique pointed lifts satisfy
\[
 \bigl|jA_{R,1}(0,R)-j(0,R)\bigr|
 \geq\frac{R}{1+\sqrt2R}>\frac25.
\]
Since $d_Q\geq d$, the same lower bound holds for $d^{\uparrow}$.
This verifies the example of LM4 C.92 directly.  Each $A_{R,t}$
is itself an allowed isotopy.  The sequence is therefore not a
noncontractible loop and does not refute Task D.2.
\end{proof}

\statement{D.7}{Proved}{Families with one compact quotient support fill}
Let $a:S^m\to C$ be a based sphere family, $m\geq1$, and suppose
that one compact set $K\Subset N$ satisfies
$a_\xi(u)=u$ whenever $\pi u\notin K$.  Then $a$ has a based
null-homotopy in $C$ with one compact quotient support, and hence
with a uniform tail.
\begin{proof}
If $N$ is compact, D.5 applies directly.  Otherwise choose a
connected compact subsurface $S\subset N$ with boundary
and $K\subset\operatorname{int}S$.  The descents restrict to a
based family in
$\operatorname{Homeo}(S\mathbin{\mathrm{rel}}\partial S)_0$:
restriction preserves $S$ because the maps are the identity off
$K$, and paths from the sphere basepoint remain in this group.
This group is contractible by
\cite[Lemma 2.1(ii)]{DE_Yagasaki2000}.  Choose its identity-fixed
contraction, extend it by the identity to $N$, and lift the
resulting isotopy starting at the identity.  The lift is the
identity outside $\pi^{-1}(S)$ and commutes with every deck
transformation.  At the given sphere it agrees with $a$, since
both are lifts with the same value at any point outside
$\pi^{-1}(S)$.  D.2 gives joint extension to the original closed
disk.  The same construction works for all inverse maps.
\end{proof}

\statement{D.8}{Proved}{The trivial operator case and the compact-open exception}
If $Q=1$, the whole group $C$ and every point stabilizer $F$ are
contractible in their actual topology, even though the identity
component of $\operatorname{Homeo}(\mathbb R^2\mathbin{\mathrm{rel}}0)$
in the interior compact-open topology has homotopy type $S^1$.
\begin{proof}
In standard disk coordinates the Alexander formula
\[
 H_t(a)(z)=
 \begin{cases}t\,a(z/t),&|z|\leq t,\quad t>0,\\
 z,&|z|\geq t\text{ or }t=0
 \end{cases}
\]
is continuous for maps and inverses, equals $a$ at $t=1$, and
equals $1$ at $t=0$.  Its displacement on the radius-$t$ disk
is at most $2t$, so continuity at $t=0$ is uniform in $a$.
It fixes $1$ and preserves the stabilizer of the origin.
Conjugation by a boundary-fixed disk homeomorphism taking a
specified interior point to the origin handles every $F$.
The compact-open statement is exactly the pointed-plane exception
in \cite[Theorem 1.1(i)]{DE_Yagasaki2000}; it is not a statement
about the stronger topology of D.3.
\end{proof}

\statement{D.9}{Open}{The arbitrary noncompact quotient uniform-tail deformation}
Assume the quotient $N$ is noncompact.  Let $P$ be a compact
metrizable space and $a:P\to F$ a continuous map.
Does $a$ admit a homotopy to the constant identity that is
continuous in the metric $d^{\uparrow}$?
The relative version for a prescribed subset on which $a=1$ is
also undecided here.  In particular we have not proved or
refuted $\pi_1(C,1)=0$ in this generality.

The precise missing estimate is \eqref{eq:DE-tail} uniformly in
$(s,t)\in P\times[0,1]$.  The unqualified compact-open contraction
of a pointed noncompact surface does not supply it, as D.6 makes
concrete.  A sufficient constructive route is a deformation of
every based sphere family to one with common compact quotient
support, preserving this estimate; D.7 would then finish the
sphere.  No such cutoff has been constructed for an arbitrary
$d_Q$.

We checked the later uniform-deformation literature as well.
Yagasaki's covering-space theorem has a compact base
\cite[Theorem 1.1]{DE_Yagasaki2014}.  Its end theorem assumes
specified Euclidean ends.  The later end theorem requires finitely
many isolated proper product ends with the stated end-deformation
property \cite[Theorem 1.3]{DE_Yagasaki2018}.  Neither hypothesis
follows from the present data.  In particular the bounded metric
$d_Q$ need not be a proper end metric, and no comparison transferring
the needed deformation to $d^{\uparrow}$ has been proved.  We also
considered equivariant deformations of a null lifted triangulation.
Its pointwise skeleton stabilizer contracts.  The family of skeleton
embeddings still requires a deformation with the same uniform control.
The escaping sequence in D.6 has not been
turned into a nontrivial loop; presenting it as one would be an
error.

\statement{D.10}{Open}{Arbitrary-frontier descent in the exact comparison}
Let $U$ be a simply connected open orbit in the original surface.
The prime-end disk need not map continuously onto its closure.
When the interior identification has such an extension,
exists, its boundary map may fail to be injective.
We have not proved that an intrinsic centralizer filling can
be chosen to extend over every such original frontier, with
jointly continuous inverses.  An interior conjugacy alone does
not suffice.  One must retain all original approach limits and
prescribed frontier values; a boundary equivalence relation
describes only the cases where a continuous boundary map exists.
This is the
arbitrary-frontier descent obligation listed in LM4 C.12.

\statement{D.11}{Open}{Nonfixed-frontier pasting and infinitely many orbit types}
The remaining LM4 C.12 comparison requires simultaneous fillings
on all open orbits agreeing with a deforming, possibly nonfixed,
nonopen locus, including infinitely many operator-inequivalent
remaining components and accumulation outside the treated closed
sets.  We have not proved or refuted this assertion.  The missing
conditions are (i) existence of mutually compatible restrictions
on every common frontier, (ii) extension of these restrictions over
the incident open orbits, and (iii) uniform control of both maps
and inverses across arbitrary accumulating families.  D.2 controls
one original disk with fixed boundary and does not supply these
three comparisons.

\statement{D.12}{Open}{The universal admissible level-zero kernel}
For an admissible punctured surface, with actual homeomorphism
group $A$, a strict action $u:\Gamma\to A$, and $H=u(\Gamma)$,
the unresolved comparison of LM4 C.12 is injectivity of
\[
 \pi_1(C_A(H),1)\longrightarrow\pi_1(A,1)^\Gamma.
\]
For hyperbolic finite-type surfaces the target is zero, so the
question is whether every based centralizer loop contracts in
the centralizer.  For the plane, annulus, and torus the respective
targets are the invariant subgroups of $\mathbb Z$, $\mathbb Z$,
and $\mathbb Z^2$.  No counterexample inside this admitted
category and no proof covering D.9--D.11 is obtained here.  The
unpunctured sphere example mentioned in LM4 C.12 is outside the
stated category and cannot settle this assertion.

\statement{D.13}{Proved}{Planar filling hulls preserve the invariant diameter bound}
Choose closed-disk coordinates in which $d$ is Euclidean, and retain
$d_Q(u,v)=\sup_q|qu-qv|$.  This changes neither the actual group
topology nor its $C_0$ tail condition.  For a nonempty compact
$A\subset O$, let $\operatorname{Fill}(A)$ be $A$ together with
all relatively compact components of $O\setminus A$.  Then
\[
 q\operatorname{Fill}(A)=\operatorname{Fill}(qA),\qquad
 \operatorname{diam}_{d_Q}\operatorname{Fill}(A)
 =\operatorname{diam}_{d_Q}(A).
\]
Moreover, if $A\subset B_{d_Q}(x,r)$ is compact, then
$\operatorname{Fill}(A)\subset B_{d_Q}(x,r)$.  Every connected
component of a $d_Q$-ball is a simply connected open planar domain.
In particular every loop contained in one such ball has a filling
in that same ball.
\begin{proof}
Two compatible metrics on the compact disk have uniform moduli of
comparison.  Taking their suprema over $Q$ preserves these moduli,
so the invariant uniform topologies and vanishing-tail conditions
are identical.

Homeomorphisms of $O$ preserve relatively compact complementary
components.  This proves the equivariance
identity.  In Euclidean disk coordinates,
\[
 \operatorname{Fill}(qA)\subset\operatorname{conv}(qA)\subset O.
\]
The first inclusion holds because a point outside the convex hull
can be joined to the exterior without meeting $qA$; the second
uses compactness of $qA$ and convexity of the open disk.  Consequently,
for $u,v\in\operatorname{Fill}(A)$ and $q\in Q$,
\[
 |qu-qv|\leq\operatorname{diam}_{d}(qA)
 \leq\operatorname{diam}_{d_Q}(A).
\]
Taking suprema gives one inequality; $A\subset\operatorname{Fill}(A)$
gives the reverse.  If $A\subset B_{d_Q}(x,r)$, compactness gives
$r'<r$ with $d_Q(a,x)\leq r'$ for all $a\in A$.  Every $qA$ lies
in the Euclidean radius-$r'$ ball about $qx$, so its convex hull
does too.  Thus $\operatorname{Fill}(A)\subset B_{d_Q}(x,r)$.

Let a loop lie in a connected component $W$ of this open ball.
Its image is compact and connected.  Its filled hull is connected
and lies in the ball, hence in $W$.  A compact full connected planar
set has arbitrarily small simply connected polygonal neighborhoods:
take a finite union of sufficiently fine grid squares meeting it,
thicken slightly, and fill the bounded complementary components.
The hull containment just proved keeps such a neighborhood in
$W$ if the grid is sufficiently fine.  The loop therefore contracts
in a polygonal disk contained in $W$.  This proves simple
connectedness.  Notice the exact limit of the conclusion: it
controls fillings of point trajectories.  It does not extend an
entire parameter family to equivariant ambient homeomorphisms,
which is the additional step required in D.9.
\end{proof}

\statement{D.14}{Proved}{An inner-diameter metric retaining every based sphere family}
With the metric in D.13, put
\[
 \rho(u,v)=\inf\{\operatorname{diam}_{d_Q}(\alpha):
   \alpha\subset O\text{ is an arc joining }u\text{ to }v\}.
\]
Then $\rho$ is a bounded $Q$-invariant metric inducing the original
interior topology, and $d_Q\leq\rho$.  Its open balls are path
connected.  If $U(u,r)$ is the component of $B_{d_Q}(u,r)$
containing $u$, then $U(u,r)$ is an open topological disk and
\[
 B_\rho(u,r)\subset U(u,r)\subset B_\rho(u,2r).
\]
Every map $S^j\to B_\rho(u,r)$, $j\geq1$, is
null-homotopic in $B_\rho(u,2r)$.

Let $F^\rho\subset F$ consist of maps whose lift displacement for
$\rho$ belongs to $C_0(N)$, equipped with the uniform $\rho$ topology
on maps and inverses.  Every based continuous family $P\to F$ from
a compact, connected, locally path-connected metric space is also
a continuous family into $F^\rho$, with a uniform $\rho$ tail.
Consequently inclusion $F^\rho\to F$ induces isomorphisms on every
positive based homotopy group.  This reduces the sphere-family
part of D.9 to a metric with connected balls and the displayed
uniform filling bound; it does not prove the ambient deformation.
\begin{proof}
The inequality $d_Q\leq\rho$ and positive definiteness are immediate.
Concatenating two arcs and extracting an arc from their path image
proves the triangle inequality, since the two diameters add when
the sets meet.  Invariance follows by applying $q$ to arcs.
Boundedness follows from boundedness of $d_Q$ on all of $O$.
Local path connectedness of the original surface implies that
two sufficiently nearby points can be joined in an arbitrarily
small $d_Q$ neighborhood; hence $\rho$ induces the same topology.
If $\rho(u,v)<r$, choose an arc from $u$ to $v$ with $d_Q$ diameter
less than $r$.  Every initial subarc has the same bound, so the
whole arc lies in $B_\rho(u,r)$.  This proves path connectedness.
The same arc lies in $U(u,r)$, giving the first displayed
inclusion.  Conversely, an arc in $U(u,r)$ from $u$ to any of
its points has $d_Q$ diameter strictly less than $2r$ by compactness.
This gives the second inclusion.  D.13 makes $U(u,r)$ simply
connected; the classification of open simply connected planar
domains makes it an open disk (including the case $U(u,r)=O$).

For a sphere map into that ball, add one arc in the ball from
$u$ to a point of its image.  The union $A$ is compact and connected,
and lies in $B_{d_Q}(u,r')$ for some $r'<r$.  By D.13 its filled
hull lies in the same ball.  Choose a simply connected polygonal
neighborhood $W$ of that hull contained in
$B_{d_Q}(u,r'')$, where $r'<r''<r$.  Any point of $W$ can be joined
to $u$ inside $W$, on an arc of $d_Q$ diameter at most $2r''$.
Thus $W\subset B_\rho(u,2r)$, and contractibility of $W$ fills the
sphere.  The choice of $W$ is needed only for this individual
sphere; no continuous selection of neighborhoods is asserted.

Every member of $C$ is uniformly continuous for $\rho$: a modulus
$\omega$ for its $d_Q$ uniform continuity bounds the diameter of
the image of every arc by $\omega$ of its diameter.  Taking infima,
with an arbitrarily small error, proves the same modulus for
$\rho$.  Apply this also to the inverse.  Thus the indicated
uniform group topology is well defined.

Let $a:P\to F$ be as in the statement.  A compact locally
path-connected metric space is uniformly locally path connected:
for every positive parameter diameter there is a neighborhood
size such that any sufficiently close pair of points can be
joined by a path of that diameter.  This follows from a finite
cover by connected open sets and a Lebesgue number.  Uniform
continuity of $a$ in the uniform $d_Q$ topology therefore implies
uniform $\rho$ continuity: along such a short parameter path,
the trajectory of any $v\in O$ has $d_Q$ diameter less than
$2\epsilon$, giving $\rho(a_sv,a_tv)<2\epsilon$.  The argument
applies to the inverse family as well.

Finally, D.2 supplies a common tail on which
$d_Q(a_sv,v)<\epsilon$ for all $s\in P$.  Join the basepoint of
$P$, where $a=1$, to any $s$.  The resulting trajectory of $v$
stays in $B_{d_Q}(v,\epsilon)$, so it has diameter at most
$2\epsilon$ and $\rho(a_sv,v)\leq2\epsilon$.  This proves the
uniform $\rho$ tail, including inverses.  Apply the result to
based spheres and to based homotopies between them.  Surjectivity
and injectivity on each positive homotopy group follow respectively.
\end{proof}
